% ----------------------
\subsection{Section 59 (7 August 1831)}
% ----------------------
	\label{DC59}
	
	% -----
	\subsubsection{Historical Background}
	% -----
		\label{DC59-HB}
		
		% --
		\paragraph{JSP}
		% --
		
			``On 7 August 1831, JS dictated a revelation in Missouri `instructing the sa[i]nts how to keep the sabath \& how to fast and pray.' The revelation was specifically addressed to those `who have come up unto' Missouri, in fulfillment of the commandment to gather there and build up the city of Zion. Some of the instruction in the revelation probably came in response to the conduct of the inhabitants of Independence, Jackson County, Missouri, among whom these Saints were living. Many of those already in Jackson County had migrated there from southern states, whereas most church members entering the area were from the Northeast. As William W. Phelps, who traveled with JS to Missouri, explained in a July 1831 letter, Jackson County residents were `emigrants from Tennessee, Kentucky, Virginia, and the Carolinas, \&c., with customs, manners, modes of living and a climate entirely different from the northerners.' One custom that was especially different was Sabbath day observance. A later JS history characterizes many of the residents as `the basest of men' who `had fled from the face of civilized society, to the frontier country to escape the hand of justice, in their midnight revels, their sabbath breaking, horseracing, and gambling.' A traveler to western Missouri in 1833 made a similar observation, stating that `the only indications of its being Sunday' in the area was `the unusual Gambling \& noise, \& assemblies around taverns.' Sabbath day observance, however, was an important component of worship to many members of the Church of Christ. Perhaps because of the general nonobservance of the Sabbath among the inhabitants of Jackson County, the 7 August revelation emphasized the importance of keeping the Sabbath day holy, outlining what church members should do on that day. These guidelines filled a void that neither the `Articles and Covenants' of the church nor the February 1831 revelation giving the `Laws of the Church of Christ' had addressed, thereby providing direction to those who would be building up the city of Zion without the benefit of JS’s in-person leadership.
			
			``The revelation may have also been a response to the concerns of those who had gone to Missouri and felt daunted by the task of building up Zion in a region described by one observer as containing only `two or three merchants stores, and fifteen or twenty dwelling houses, built mostly of logs hewed on both sides.' The writer Washington Irving, who traveled through Independence in 1832 on an expedition with federal Indian commissioners, also commented on the `rougher and rougher life' he encountered as he got closer to the town, while his traveling companion Charles Latrobe described Independence as `full of promise' but containing `nothing but a ragged congeries of five or six rough log huts, two or three clapboard houses, two or three so-called hotels, alias grogshops; [and] a few stores.' Perhaps to encourage the Saints in such conditions, the revelation promised the bounties of the earth to church members and reminded them to express gratitude to God.
			
			``The revelation assured heavenly rewards for the obedient who would die in Zion---prompted, perhaps, by the death on the morning of 7 August of Polly Peck Knight, the fifty-seven-year-old wife of Joseph Knight Sr., and a friend of JS and his family. Polly Knight had traveled to Missouri with the Colesville Saints, but after falling ill she became `the first death in the church in this land.'%
				\footnote{%
					% \label{}%
					HDDC 752 contains an account of Polly Knight's death.
					}
			It is unclear whether this revelation was dictated before or after JS was informed of her death.
			
			``Oliver Cowdery served as the scribe for the original inscription of this revelation. The copy featured here belonged to Newel K. Whitney and is also in Cowdery's handwriting. Whitney's copy may be the original but is more likely a fair copy. It was likely made for him sometime after Cowdery returned to Ohio at the end of August. Around that same time, John Whitmer copied the revelation into Revelation Book 1. That there are few differences between the two copies suggests they were made around the same time or from the same copy.''
		
	% -----
	\subsubsection{Versions}
	% -----
		\label{DC59-V}
		
		See the \href{https://drive.google.com/file/d/1goAvNz8jS4R8slYfMG_db55Ty2z_vmgK/view?usp=sharing}{sections comparison} document.
	
		\begin{compactdesc}
			\item[JSP] \hyperref[EMS-1831-JS-7-AUG-1831]{7 August 1831}
			\item[EMS] 1:2 (July 1832), 9
			\item[BC] Chapter 60, 140--142
			\item[DC 1835] Section 19, 140--141
			\item[TS] 5:5 (1 March 1844), 450--451
			\item[DC 1844] Section 19, 200--202
			\item[MS] 5:9 (February 1845), 132--133
		\end{compactdesc}

	% -----
	\subsubsection{Section 59:1} % *DC59-1 
	% -----
		\label{DC59-1}

		BEHOLD, blessed, saith the Lord, are they who have come up unto this land with an eye \hyperref[SINGLE]{single} to my glory, according to my commandments.

		% \begin{framed}
		% \scriptsize
			% \begin{compactdesc}
				% \item[JSP] 
				% % \item[BC] 
				% % \item[]
			% \end{compactdesc}
		% \end{framed}

	% -----
	\subsubsection{Section 59:2} % *DC59-2 
	% -----
		\label{DC59-2}

		For those that live shall inherit the earth, and those that die shall rest from all their labors, and their works shall follow them; and they shall receive a crown in the mansions of my Father, which I have prepared for them.

		% \begin{framed}
		% \scriptsize
			% \begin{compactdesc}
				% \item[JSP] 
				% % \item[BC] 
				% % \item[]
			% \end{compactdesc}
		% \end{framed}

	% -----
	\subsubsection{Section 59:3} % *DC59-3 
	% -----
		\label{DC59-3}

		Yea, blessed are they whose feet stand upon the land of Zion, who have obeyed my gospel; for they shall receive for their reward the good things of the earth, and it shall bring forth in its strength.

		% \begin{framed}
		% \scriptsize
			% \begin{compactdesc}
				% \item[JSP] 
				% % \item[BC] 
				% % \item[]
			% \end{compactdesc}
		% \end{framed}

	% -----
	\subsubsection{Section 59:4} % *DC59-4 
	% -----
		\label{DC59-4}

		And they shall also be crowned with blessings from above, yea, and with commandments not a few, and with revelations in their time---they that are faithful and diligent before me.

		% \begin{framed}
		% \scriptsize
			% \begin{compactdesc}
				% \item[JSP] 
				% % \item[BC] 
				% % \item[]
			% \end{compactdesc}
		% \end{framed}

	% -----
	\subsubsection{Section 59:5} % *DC59-5 
	% -----
		\label{DC59-5}

		Wherefore, I give unto them a commandment, saying thus: Thou shalt love the Lord thy God with all thy heart, with all thy might, mind, and strength;%
			\footnote{%
				% \label{}%
				See \hyperref[DC4]{DC 4}.
				}
		and in the name of Jesus Christ thou shalt serve him.

		% \begin{framed}
		% \scriptsize
			% \begin{compactdesc}
				% \item[JSP] 
				% % \item[BC] 
				% % \item[]
			% \end{compactdesc}
		% \end{framed}

	% -----
	\subsubsection{Section 59:6} % *DC59-6 
	% -----
		\label{DC59-6}

		Thou shalt love thy neighbor as thyself.  Thou shalt not steal; neither commit adultery, nor kill, nor do anything like unto it.

		% \begin{framed}
		% \scriptsize
			% \begin{compactdesc}
				% \item[JSP] 
				% % \item[BC] 
				% % \item[]
			% \end{compactdesc}
		% \end{framed}

	% -----
	\subsubsection{Section 59:7} % *DC59-7 
	% -----
		\label{DC59-7}

		Thou shalt thank the Lord thy God in all things.%
		\footnote{%
			% \label{}%
			Gratitude as a commandment here.
			} 

		% \begin{framed}
		% \scriptsize
			% \begin{compactdesc}
				% \item[JSP] 
				% % \item[BC] 
				% % \item[]
			% \end{compactdesc}
		% \end{framed}

	% -----
	\subsubsection{Section 59:8} % *DC59-8 
	% -----
		\label{DC59-8}

		\footnote{%
			% \label{}%
			See \hyperref[3NEPHI-9-20]{3 Nephi 9:20} and notes there for this sacrifice.
			}%
		Thou shalt offer a sacrifice unto the Lord thy God in righteousness, even that of a broken heart and a \hyperref[CONTRITE]{contrite} spirit.

		% \begin{framed}
		% \scriptsize
			% \begin{compactdesc}
				% \item[JSP] 
				% % \item[BC] 
				% % \item[]
			% \end{compactdesc}
		% \end{framed}

	% -----
	\subsubsection{Section 59:9} % *DC59-9 
	% -----
		\label{DC59-9}

		And that thou mayest more fully keep thyself unspotted%
		\footnote{%
			% \label{}%
			``Unspotted'' occurs twice in the scriptures: here, and at \hyperref[JAMES-1-27]{James 1:27}.
			}
		from the world, thou shalt go to the house of prayer and offer up thy sacraments upon my holy day;

		% \begin{framed}
		% \scriptsize
			% \begin{compactdesc}
				% \item[JSP] 
				% % \item[BC] 
				% % \item[]
			% \end{compactdesc}
		% \end{framed}

	% -----
	\subsubsection{Section 59:10} % *DC59-10 
	% -----
		\label{DC59-10}

		For verily this is a day appointed unto you to rest from your labors, and to pay thy devotions unto the Most High;

		% \begin{framed}
		% \scriptsize
			% \begin{compactdesc}
				% \item[JSP] 
				% % \item[BC] 
				% % \item[]
			% \end{compactdesc}
		% \end{framed}

	% -----
	\subsubsection{Section 59:11} % *DC59-11 
	% -----
		\label{DC59-11}

		Nevertheless thy vows shall be offered up in righteousness on all days and at all times;

		% \begin{framed}
		% \scriptsize
			% \begin{compactdesc}
				% \item[JSP] 
				% % \item[BC] 
				% % \item[]
			% \end{compactdesc}
		% \end{framed}

	% -----
	\subsubsection{Section 59:12} % *DC59-12 
	% -----
		\label{DC59-12}

		But remember that on this, the Lord's day, thou shalt offer thine oblations%
			\footnote{%
				% \label{}%
				The 1828 dictionary only has the following definition for ``oblation'': ``Any thing offered or presented in worship or sacred service; an offering; a sacrifice.''
				}
		and thy sacraments unto the Most High, confessing thy sins unto thy brethren, and before the Lord.

		% \begin{framed}
		% \scriptsize
			% \begin{compactdesc}
				% \item[JSP] 
				% % \item[BC] 
				% % \item[]
			% \end{compactdesc}
		% \end{framed}

	% -----
	\subsubsection{Section 59:13} % *DC59-13 
	% -----
		\label{DC59-13}

		And on this day thou shalt do none other thing, only let thy food be prepared with \hyperref[SINGLE]{singleness} of heart that thy fasting may be perfect, or, in other words, that thy joy may be full.

		% \begin{framed}
		% \scriptsize
			% \begin{compactdesc}
				% \item[JSP] 
				% % \item[BC] 
				% % \item[]
			% \end{compactdesc}
		% \end{framed}

	% -----
	\subsubsection{Section 59:14} % *DC59-14 
	% -----
		\label{DC59-14}

		Verily, this is fasting and prayer, or in other words, rejoicing and prayer.

		% \begin{framed}
		% \scriptsize
			% \begin{compactdesc}
				% \item[JSP] 
				% % \item[BC] 
				% % \item[]
			% \end{compactdesc}
		% \end{framed}

	% -----
	\subsubsection{Section 59:15} % *DC59-15 
	% -----
		\label{DC59-15}

		And inasmuch as ye do these things with thanksgiving, with cheerful hearts and countenances, not with much laughter, for this is sin, but with a glad heart and a cheerful countenance---

		% \begin{framed}
		% \scriptsize
			% \begin{compactdesc}
				% \item[JSP] 
				% % \item[BC] 
				% % \item[]
			% \end{compactdesc}
		% \end{framed}

	% -----
	\subsubsection{Section 59:16} % *DC59-16 
	% -----
		\label{DC59-16}

		Verily I say, that inasmuch as ye do this, the fulness of the earth is yours, the beasts of the field and the fowls of the air, and that which climbeth upon the trees and walketh upon the earth;

		% \begin{framed}
		% \scriptsize
			% \begin{compactdesc}
				% \item[JSP] 
				% % \item[BC] 
				% % \item[]
			% \end{compactdesc}
		% \end{framed}

	% -----
	\subsubsection{Section 59:17} % *DC59-17 
	% -----
		\label{DC59-17}

		Yea, and the herb, and the good things which come of the earth, whether for food or for raiment, or for houses, or for barns, or for orchards, or for gardens, or for vineyards;

		% \begin{framed}
		% \scriptsize
			% \begin{compactdesc}
				% \item[JSP] 
				% % \item[BC] 
				% % \item[]
			% \end{compactdesc}
		% \end{framed}

	% -----
	\subsubsection{Section 59:18} % *DC59-18 
	% -----
		\label{DC59-18}

		Yea, all things which come of the earth, in the season thereof, are made for the benefit and the use of man, both to please the eye and to gladden the heart;

		% \begin{framed}
		% \scriptsize
			% \begin{compactdesc}
				% \item[JSP] 
				% % \item[BC] 
				% % \item[]
			% \end{compactdesc}
		% \end{framed}

	% -----
	\subsubsection{Section 59:19} % *DC59-19 
	% -----
		\label{DC59-19}

		Yea, for food and for raiment, for taste and for smell, to strengthen the body and to enliven the soul.

		% \begin{framed}
		% \scriptsize
			% \begin{compactdesc}
				% \item[JSP] 
				% % \item[BC] 
				% % \item[]
			% \end{compactdesc}
		% \end{framed}

	% -----
	\subsubsection{Section 59:20} % *DC59-20 
	% -----
		\label{DC59-20}

		And it pleaseth God that he hath given all these things unto man; for unto this end were they made to be used, with judgment, not to excess, neither by extortion.

		% \begin{framed}
		% \scriptsize
			% \begin{compactdesc}
				% \item[JSP] 
				% % \item[BC] 
				% % \item[]
			% \end{compactdesc}
		% \end{framed}

	% -----
	\subsubsection{Section 59:21} % *DC59-21 
	% -----
		\label{DC59-21}

		And in nothing doth man \hyperref[OFFENSE]{offend} God, or against none is his wrath kindled, save those who confess not his hand in all things, and obey not his commandments.

		% \begin{framed}
		% \scriptsize
			% \begin{compactdesc}
				% \item[JSP] 
				% % \item[BC] 
				% % \item[]
			% \end{compactdesc}
		% \end{framed}

	% -----
	\subsubsection{Section 59:22} % *DC59-22 
	% -----
		\label{DC59-22}

		Behold, this is according to the law and the prophets; wherefore, trouble me no more concerning this matter.

		% \begin{framed}
		% \scriptsize
			% \begin{compactdesc}
				% \item[JSP] 
				% % \item[BC] 
				% % \item[]
			% \end{compactdesc}
		% \end{framed}

	% -----
	\subsubsection{Section 59:23} % *DC59-23 
	% -----
		\label{DC59-23}

		But learn that he who doeth the works of \hyperref[RIGHTEOUSNESS]{righteousness} shall receive his reward, even peace in this world and eternal life in the world to come.

		% \begin{framed}
		% \scriptsize
			% \begin{compactdesc}
				% \item[JSP] 
				% % \item[BC] 
				% % \item[]
			% \end{compactdesc}
		% \end{framed}

	% -----
	\subsubsection{Section 59:24} % *DC59-24 
	% -----
		\label{DC59-24}

		I, the Lord, have spoken it, and the Spirit beareth record.  Amen.

		% \begin{framed}
		% \scriptsize
			% \begin{compactdesc}
				% \item[JSP] 
				% % \item[BC] 
				% % \item[]
			% \end{compactdesc}
		% \end{framed}